\documentclass[11pt]{article}
\usepackage{manual_docs_style}

% Source-review destinations used by the docs comment synchronizer.
\DeclareRobustCommand{\ReviewAnchor}[2]{%
  \ifcsname hypertarget\endcsname
    \hypertarget{review-anchor:\detokenize{#1}}{}%
  \else
    \pdfdest name {review-anchor:#1} xyz%
  \fi
}
% Keep review anchors out of PDF bookmark strings.
\ifcsname pdfstringdefDisableCommands\endcsname
  \pdfstringdefDisableCommands{%
    \def\ReviewAnchor#1#2{}%
  }
\fi
\begin{document}

\docTitle{Terminology and Common Utilities}
\phantomsection\label{docstart:terminology_common_utilities}

\agentProposedEdit
  {agent-add-b52bdc65-df6d-57e9-b82f-0e13e43012d4}
  {Add this agent-authored block for review.}
  {}
  {\textbf{Paper-aligned terminology.}}
\begin{itemize}
\item A \name{} task is a batched closed-set root-cause attribution problem in which the agent assigns one class to every test sample.
\item A time-series sample is one multivariate observation with \(T\) time steps and \(C\) observed channels.
\item The test batch is the unlabeled set of samples for which the agent must return predictions.
\item The optional labeled support batch contains samples from the same simulator as the test batch together with their ground-truth labels.
\item An evaluation episode is one complete agent interaction with an instantiated task, beginning when the agent receives the task and ending when the agent returns control.
\item A physical simulator generates a noise-free trajectory from a parameter vector and an initial state.
\item An intervention sample contains one instantaneous step change in one candidate simulator parameter at an intervention time.
\item A no-intervention sample keeps every simulator parameter fixed throughout the simulation and has class label \(\mathrm{NI}\).
\item The observation model \(p_{\kappa}\) generates a noisy time-series observation from a noise-free trajectory, with the scalar \(\kappa\) controlling the intensity of measurement corruption.
\end{itemize}

\agentProposedEdit
  {agent-add-233fa2d6-8adb-5264-9381-e8f21aca7f23}
  {Add this agent-authored block for review.}
  {}
  {\textbf{Paper-to-implementation mappings.}}
\begin{itemize}
\item The paper's with-domain-context and without-domain-context settings map to \code{desc_level=high} and \code{desc_level=none}, respectively.
\item The paper's \(\kappa_{\mathrm{low}}\) and \(\kappa_{\mathrm{high}}\) observation-noise regimes map to \code{noise_level=low} and \code{noise_level=high}, respectively.
\item The internal \code{noise_level=none} setting materializes a noise-free trajectory and is not one of the paper's evaluated low-noise or high-noise regimes.
\item The paper's labeled-examples axis maps to the internal support-sample count \code{shot} and its resolved grouping key \code{shot_slug}.
\item The paper's direct-answer mode \(\mathcal{S}_{\mathrm{direct}}\) and programmatic-submission mode \(\mathcal{S}_{\mathrm{prog}}\) map to \code{task_mode=direct} and \code{task_mode=code}, respectively.
\item Internal generation identifiers such as \code{simulator_recipe}, \code{run_id}, and \code{family_id} belong in the corresponding pipeline-stage documentation rather than in the paper terminology glossary.
\end{itemize}

\section*{\agentProposedEdit{agent-remove-ed0fa7b2-2aa3-515b-ace0-257adad89635}{replaced by the paper-aligned glossary and mappings above}{Documentation Terminology}{}}
\begin{itemize}
\item \agentProposedEdit{agent-remove-e44b6e6d-c161-50f9-b3d5-0ee4a94572d1}{minimum absolute distance is an internal configuration field and is not part of the paper terminology glossary}{(\termMinAbsDist): minimum absolute distance/change required between values.}{}
\agentProposedEdit{agent-remove-81f066c5-6755-5d41-9e7c-5640857e99ab}{the allowed simulator registry belongs in the environment documentation}{The allowed values are defined in \code{ALLOWED_TRACEBENCH_MODELS}.}{}
\item \agentProposedEdit{agent-remove-00a4d24f-06f2-5bd0-8dd4-1ab4197e0235}{simulator_recipe is an internal generation identifier that belongs in the pipeline-stage documentation}{(\termSimulatorRecipe): One configuration that can be given as input to the simulate stage to generate.}{}
\item \agentProposedEdit{agent-remove-11f967a1-522b-5309-983b-deb3eef71606}{run_id is an internal generation identifier that belongs in the pipeline-stage documentation}{(\termRunID): An identifier for a single run.}{}
\item \agentProposedEdit{agent-remove-098c62e1-6db8-5ed1-b49f-4ea3bfe012c6}{family is an internal generation concept that belongs in the pipeline-stage documentation}{(\termFamily): A \termFamily is the baseline-centered group of run nodes generated from one sampled baseline recipe.}{}
\agentProposedEdit{agent-remove-72fb93c8-7739-5635-9022-e4213528e949}{family is an internal generation concept that belongs in the pipeline-stage documentation}{It contains the baseline run, all intervention children derived from that baseline, and the matched time-zero/static-change probe for each intervention.}{}
\agentProposedEdit{agent-remove-bfcbf6ed-c22a-5528-a139-f16b6344cae3}{family_id is a separate baseline-recipe hash rather than the baseline run_id}{All nodes in the group share the same \code{family_id}, which is \termRunID of the baseline.}{}
\item \agentProposedEdit{agent-remove-5225c385-9c63-58d1-bd8b-8b3a7b8f81c1}{replaced by the paper-aligned observation-noise definition and mapping above}{(\termNoiseLevel): The observation-noise profile applied when materializing a sample.}{}
\end{itemize}


\section*{\agentProposedEdit{agent-remove-a7391171-2fc1-550e-9d40-bca601757405}{repository layout is operational documentation rather than paper terminology}{Private repository structure}{}}
\agentProposedEdit
  {agent-add-db6cc8d4-9968-53e1-b90e-83421718b2fd}
  {Add this agent-authored block for review.}
  {}
  {The complete private-repository listing below should be removed from this terminology page and maintained only in repository-specific operational documentation.}
\begin{DocCode}
pyproject.toml      # project metadata + dependencies1
uv.lock             # exact locked dependency versions
TraceBench/
models/simulink/
workflows/
terminal-bench/
simscape_custom/
\end{DocCode}

\section*{\agentProposedEdit{agent-remove-f4df8f6f-e678-53a3-8b3e-1afc75fbab6b}{machine ownership is operational documentation rather than paper terminology}{Computer and responsibily}{}}
\agentProposedEdit
  {agent-add-70b54e91-11f0-5634-8e37-ec55c9681702}
  {Add this agent-authored block for review.}
  {}
  {The complete machine-responsibility section below should be removed from this terminology page and maintained only in operational documentation.}
\begin{itemize}
    \item \code{MacBook}: documentation and paper writing.
    \item \code{L7161}: generate large amounts of data.
    \item \code{tommaso-desktop}: experiment running.
    \item \code{sip-ml}: student work.
\end{itemize}



\section*{\agentProposedEdit{agent-remove-bf6b95b8-454c-589e-90eb-a333a4cb1107}{utility implementation details belong in the relevant pipeline-stage documentation}{Common Utilities}{}}
\agentProposedEdit
  {agent-add-1dc6f19a-fc5a-5043-9473-94482bc05147}
  {Add this agent-authored block for review.}
  {}
  {The complete utility-pseudocode section below should be removed because it documents implementation details rather than paper terminology.}

\begin{algorithm}[H]
  \caption{\code{materialize(uuid_path, noise_level="none", seed, uuid_reference_path=None) -> (df, noise_analysis)}}
\label{def:materialize-utility}
\textbf{Purpose.} Loads a saved run artifact, optionally applies the requested observation-noise profile, anonymizes columns to \code{col1}, \code{col2}, \ldots, \code{colN}, and returns the agent-visible dataframe.

\textbf{Used by.} Cheap Filter Metrics, Agentic Rollout, and Final Artifact Scoring.
\par\smallskip
\begin{algorithmic}[1]
\State \code{src} $\gets$ read parquet file at \code{uuid_path}
\State \code{ref} $\gets$ \code{None}
\State \code{noise_analysis} $\gets$ \code{None}
\If{\code{uuid_reference_path} is not \code{None}}
    \State \code{ref} $\gets$ read parquet file at \code{uuid_reference_path}
\EndIf
\If{\code{noise_level} $\neq$ \code{"none"}}
    \State \code{(noisy_df, noise_analysis)} $\gets$ \code{add_noise(src, seed, noise_level, ref=ref)}
    \State \code{df} $\gets$ \code{noisy_df}
\Else
    \State \code{df} $\gets$ \code{src}
\EndIf
\State Rename all columns of \code{df} to \code{col1}, \code{col2}, \ldots, \code{colN}
\Assert \code{colN} is time (i.e. regularly sampled)
\State \Return \code{(df, noise_analysis)}
\end{algorithmic}
\end{algorithm}

Where \code{add_noise} is defined for each \termSimulator in \code{noise_adder.py}.

\agentProposedEdit{agent-remove-cd5db0d6-d948-5ecb-b39c-2e48383e0f6d}{NOISE_DICT now contains separate low and high profile objects rather than low-only base magnitudes}{Each \termSimulator's \code{noise_adder.py} defines \code{NOISE_DICT}, whose values give the \termSimulator-specific base magnitudes for the low-noise regime.}{}
\begin{algorithm}[H]
\caption{\code{add_noise(...)}}
\label{def:add-noise-utility}
\textbf{Purpose.} Applies the \termSimulator-specific low or high observation-noise profile to a clean simulator trajectory and optionally to its matched reference trajectory.

\textbf{Used by.} \code{materialize(...)}, Cheap Filter Metrics, Agentic Rollout, and noise-analysis reporting.
\par\smallskip
\begin{algorithmic}[1]
\Assert{\code{noise_level} is in \code{"none","low","high"}}
\State \code{current_seed} $\gets$ \code{seed}
\State \code{scale} $\gets$ \code{NOISE_LEVEL_MULTIPLIER[noise_level]}
\State \code{src_noisy} $\gets$ apply \termSimulator-specific noise to \code{src} using \code{current_seed}, \code{scale}, \code{NOISE_DICT}
\If{\code{ref} is not \code{None}}
    \State \code{ref_noisy} $\gets$ apply \termSimulator-specific noise to \code{ref} using \code{current_seed+10}, \code{scale}, \code{NOISE_DICT}
\Else
    \State \code{ref_noisy} $\gets$ \code{None}
\EndIf
\State \code{noise_analysis} $\gets$ \code{quantify_effect_snr(src, src_noisy, ref, ref_noisy)}
\State \Return \code{(src_noisy, noise_analysis)}

\end{algorithmic}
\end{algorithm}

\begin{algorithm}[H]
\caption{\code{quantify_effect_snr(...)}}
\label{def:quantify-noise-utility}
\textbf{Purpose.} Computes the effect-to-noise signal-to-noise ratio between a clean trajectory, its noisy version, and optionally a matched reference trajectory.

\textbf{Used by.} Noise analysis during materialization and cheap-filter metrics.
\par\smallskip
\begin{algorithmic}[1]
\State \code{snr_analysis} $\gets$ \code{[]}
\ForAll{observed non-time signals \code{signal}}
    \If{\code{ref} is not \code{None}}
        \State \code{clean_effect_signal} $\gets$ \code{src[signal]} $-$ \code{ref[signal]}
        \If{\code{ref_noisy} is not \code{None}}
            \State \code{src_noise_signal} $\gets$ \code{src_noisy[signal]} $-$ \code{src[signal]}
            \State \code{ref_noise_signal} $\gets$ \code{ref_noisy[signal]} $-$ \code{ref[signal]}
            \State \code{noise_signal} $\gets$ \code{src_noise_signal} $-$ \code{ref_noise_signal}
        \Else
            \State \code{noise_signal} $\gets$ \code{src_noisy[signal]} $-$ \code{src[signal]}
        \EndIf
    \Else
        \State \code{clean_effect_signal} $\gets$ \code{src[signal]}
        \State \code{noise_signal} $\gets$ \code{src_noisy[signal]} $-$ \code{src[signal]}
    \EndIf
    \State append \code{SNR_db(clean_effect_signal, noise_signal)} to \code{snr_analysis}
\EndFor
\State \Return \code{snr_analysis}
\end{algorithmic}
\end{algorithm}

Where 
\begin{DocCode}
    SNR_db(effect_signal, noise) = 20 log10(RMS(effect_signal) / RMS(noise))
\end{DocCode}
\begin{itemize}
    \item \code{effect_signal} is the intervention effect relative to the reference trajectory, not the raw signal amplitude.
    \item Signals undergo preprocessing before the formula is applied, following \code{docs/overleaf_paper/src/appendix/appendix_environment.tex::*::Signal preprocessing}.
    \item \code{quantify_effect_snr} also accepts \textbf{list} inputs for \code{src}, \code{src_noisy}, \code{ref=None}, and \code{ref_noisy=None}.
    \item For list inputs, it returns the mean SNR by first averaging the linear effect-to-noise ratios over the list.
    \item The returned SNR in dB is then computed as \code{20*log10(mean effect-to-noise ratio)}.
    \end{itemize}
\end{document}
